\title{Subject focus in French and Spanish}
\author{Alessia Cassarà}
\BackBody{This book investigates the realization of subject focus in French and Spanish from a comparative, variationist perspective. Building on spontaneous dialogical data, it examines the range of syntactic, prosodic, and elliptical strategies speakers use to encode narrowly focused subjects. The study draws on two highly comparable sub-corpora of spontaneous speech elicited through the same task, allowing for fine-grained cross-linguistic comparison. Adopting a Question Under Discussion (QUD) framework, it analyzes how discourse factors such as focus type, givenness, and argument structure condition speakers’ choices among competing focus-marking strategies. Particular attention is paid to the alternation between full forms (e.g. clefts, postverbal subjects) and more economical elliptical answers (e.g. fragments, reduced clefts). Quantitative analyses show that subject focus realization is not free but systematically constrained by pragmatic and semantic factors, albeit differently across the two languages. The results challenge rigid form–function mappings and highlight the role of language economy in discourse. Overall, the book contributes to a more flexible typology of focus realization and to our understanding of syntactic variation in closely related Romance languages.}
\renewcommand{\lsSeries}{orl}
\renewcommand{\lsSeriesNumber}{9}
\renewcommand{\lsID}{581}
\lsCoverTitleSizes{45pt}{15mm}% Font setting for the title page
\renewcommand{\lsISBNdigital}{978-3-96110-572-4}
\renewcommand{\lsISBNhardcover}{978-3-98554-193-5}
\BookDOI{10.5281/zenodo.19564874}
\typesetter{Sebastian Nordhoff}
\proofreader{Amir Ghorbanpour,
Amy Amoakuh,
Anna Staňková,
Brett Reynolds,
David Carrasco Coquillat,
Frederic Blum,
George Walkden,
Gerald Delahunty,
Jean Nitzke,
Lachlan Mackenzie,
Mahamane Laoualy Abdoulaye,
Mary Ann Walter,
Matthew Korte,
Max Rogan,
Natalia Cáceres,
Nicoletta Romeo,
Rui Yamawaki
}

